\section{Entorno de pruebas}
\label{sec:entorno_pruebas}

Todas las pruebas se realizaron en el mismo equipo físico, comparando dos escenarios
de despliegue: Docker~Compose y Kubernetes con Minikube. La señal de entrada fue
generada por FFmpeg~\cite{ffmpeg_project}: cuatro fuentes a 1920$\times$1080@25\,fps
en formato RAW~I420, sin compresión, circulando íntegramente por la red interna del
equipo. La tabla~\ref{tab:hardware} recoge las especificaciones relevantes para
interpretar los resultados.

\begin{table}[H]
  \centering
  \caption{Especificaciones del entorno de pruebas.}
  \label{tab:hardware}
  \small
  \begin{tabular}{|l|l|}
    \hline
    \textbf{Componente} & \textbf{Especificación} \\
    \hline
    CPU            & Intel Core i9-10900X @ 3{,}70\,GHz (10 núcleos / 20 hilos) \\
    RAM            & 128\,GB DDR4 \\
    Sistema operativo & Ubuntu 22.04.5 LTS (kernel 6.8.0-111) \\
    Docker Compose & Engine 29.1.3 / Compose v2.37.1 \\
    Kubernetes     & Minikube v1.38.1 / kubectl v1.36.0 \\
    Red interna    & Loopback (sin interfaz física) \\
    \hline
  \end{tabular}
\end{table}

\section{Validación funcional}
\label{sec:validacion_funcional}

Antes de entrar en el análisis cuantitativo de rendimiento, esta sección verifica
que cada módulo desarrollado en los capítulos~3 y~4 funciona correctamente en
condiciones reales de operación.
Para cada módulo se realizaron pruebas manuales sobre el escenario Docker~Compose
con señal activa: se ejecutó la acción correspondiente desde la interfaz de control
y se comprobó que el sistema respondía de acuerdo con su especificación.
La tabla~\ref{tab:validacion_funcional} resume los resultados.

\begin{table}[H]
  \centering
  \caption{Validación funcional por módulo.}
  \label{tab:validacion_funcional}
  \small
  \begin{tabular}{|l|p{6.8cm}|p{3.6cm}|}
    \hline
    \textbf{Módulo} & \textbf{Función verificada} & \textbf{Resultado} \\
    \hline
    Composites
      & Combinación visual de dos fuentes: pantalla completa, \ac{PIP}, lateral y ponencia.
      & OK — 4 modos + mirror \\
    \hline
    Transiciones
      & Interpolación animada entre composites mediante el elemento \texttt{compositor} de GStreamer.
      & OK — sin corte visible \\
    \hline
    Overlay dinámico
      & Rótulo inferior con texto configurable y animaciones de entrada y salida.
      & OK — auto-off en cada corte \\
    \hline
    Stream blanker
      & Control del estado de emisión: LIVE, PAUSE y NOSTREAM.
      & OK — 3 estados conmutables \\
    \hline
    Audio Follows Video
      & Fundido de audio automático al cambiar la fuente de vídeo activa.
      & OK — sin intervención manual \\
    \hline
    Telemetría STATE
      & Evento periódico cada 5\,s con el estado completo del mezclador publicado en RabbitMQ.
      & OK — latencia $<$\,50\,ms \\
    \hline
    Telemetría CHANGE
      & Evento inmediato ante cualquier acción del operador sobre la interfaz.
      & OK — latencia $<$\,50\,ms \\
    \hline
    Despliegue Docker Compose
      & Sistema completo operativo con todos los servicios en estado \textit{healthy}.
      & OK — arranque $\approx$\,20\,s \\
    \hline
    Despliegue Kubernetes
      & Clúster Minikube con todos los pods en estado \textit{Ready}.
      & OK — arranque $\approx$\,90\,s \\
    \hline
  \end{tabular}
\end{table}

Todos los módulos superaron la verificación.
